\chapter{Capacitors and Capacitive Reactance}
\label{ch:capacitors}

\section{The Capacitor}

A \textbf{capacitor} stores energy in an electric field between two
conductive plates separated by an insulating \textbf{dielectric}. Its
defining property is \textbf{capacitance} $C$, measured in
\textbf{farads (F)}:
\[
\boxed{C = \frac{Q}{U}}
\]
where $Q$ is the charge stored (coulombs) and $U$ is the voltage across
the plates. One farad is a very large unit; practical capacitors are
usually specified in microfarads ($\mu$F), nanofarads (nF), or
picofarads (pF).

For a simple parallel-plate capacitor with plate area $A$, separation
$d$, and a dielectric of permittivity $\varepsilon = \varepsilon_r
\varepsilon_0$,
\[
C = \frac{\varepsilon_r \varepsilon_0 A}{d}
\]
where $\varepsilon_0 = 8.854\times10^{-12}\,\text{F/m}$ is the
permittivity of free space and $\varepsilon_r$ is the dielectric's
\textbf{relative permittivity} (air $\approx 1$, ceramic 5--10{,}000
depending on formulation, mica $\approx 5$--7). This equation explains
capacitor construction directly: more plate area or a thinner, higher-permittivity
dielectric gives more capacitance in the same physical volume, which is
why high-value capacitors use thin rolled or multilayer construction.

\section{Charging and Discharging: the Capacitor's Defining Behavior}

A capacitor opposes \emph{changes} in voltage across it, just as an
inductor (Chapter~\ref{ch:inductors}) opposes changes in current. The
current into a capacitor is proportional to how fast its voltage is
changing:
\[
i(t) = C\,\frac{dU}{dt}.
\]
Two consequences follow directly, and both are exam favorites:
\begin{itemize}
\item A capacitor blocks DC in steady state: once fully charged
($dU/dt = 0$), no current flows.
\item A capacitor initially looks like a short circuit to a sudden
voltage step (large $dU/dt$ means large initial current), then charges
toward the applied voltage.
\end{itemize}
The time-domain charge/discharge behavior (the \textbf{RC time
constant}) is developed fully in Chapter~\ref{ch:rc-rl}.

\section{Energy Stored in a Capacitor}

The energy stored in a charged capacitor is
\[
W = \frac{1}{2}CU^2 = \frac{1}{2}QU = \frac{Q^2}{2C}.
\]
This stored energy is why capacitors are used for power-supply filtering
(smoothing ripple by supplying current between rectifier pulses),
timing circuits, and as the energy-storage element in an LC tank circuit
(Chapter~\ref{ch:rlc-resonance}).

\section{Capacitors in Series and Parallel}

Unlike resistors, capacitor combination rules are ``reversed'':

\begin{keyidea}
\textbf{Parallel capacitors add directly} (more plate area, effectively):
\[
C_{\text{total}} = C_1 + C_2 + C_3 + \cdots
\]
\textbf{Series capacitors combine by the reciprocal sum} (effectively a
thicker dielectric gap):
\[
\frac{1}{C_{\text{total}}} = \frac{1}{C_1} + \frac{1}{C_2} + \frac{1}{C_3} + \cdots
\]
For exactly two capacitors in series, $C_{\text{total}} = \dfrac{C_1
C_2}{C_1+C_2}$ (product over sum, same shortcut form as parallel
resistors).
\end{keyidea}

\begin{examplebox}[Series capacitors]
Find the total capacitance of a 100\,pF and a 220\,pF capacitor in
series.

\emph{Solution.}
\[
C_{\text{total}} = \frac{100 \times 220}{100+220} = \frac{22000}{320} \approx 68.8\,\text{pF}.
\]
As with parallel resistors, the series combination is always smaller than
the smallest individual value.
\end{examplebox}

\section{Capacitive Reactance}
\label{sec:xc}

When an AC voltage is applied to a capacitor, the current leads the
voltage: current is largest exactly when voltage is changing fastest
(at the zero-crossing) and zero when voltage is at its peak (momentarily
not changing). This $90\degree$ phase relationship, combined with the
frequency-dependent \emph{magnitude} of opposition to current flow, is
called \textbf{capacitive reactance}, $X_C$, measured in ohms:
\[
\boxed{X_C = \frac{1}{2\pi f C} = \frac{1}{\omega C}}
\]

\begin{keyidea}[Reactance vs.\ resistance]
Reactance, like resistance, is measured in ohms and limits current for a
given voltage. Unlike resistance, reactance depends on frequency and does
not dissipate energy as heat --- it stores and returns energy each cycle.
Capacitive reactance \emph{decreases} as frequency increases (a
capacitor looks increasingly like a short circuit at high frequency) and
\emph{increases} without bound as frequency approaches zero (a capacitor
looks like an open circuit at DC, consistent with Section~\ref{ch:capacitors}
above).
\end{keyidea}

\begin{examplebox}[Calculating reactance]
Find the reactance of a 100\,pF capacitor at 14.2\,MHz.

\emph{Solution.}
\[
X_C = \frac{1}{2\pi f C} = \frac{1}{2\pi \times 14.2\times10^{6} \times 100\times10^{-12}} \approx 112\,\Omega.
\]
\end{examplebox}

Because current leads voltage by exactly $90\degree$ in an ideal
capacitor, the impedance of a capacitor is written as a purely imaginary
number in the phasor (complex) notation introduced in
Chapter~\ref{ch:ac-theory}:
\[
Z_C = -jX_C = \frac{1}{j\omega C}.
\]
This notation becomes essential once capacitors and resistors are
combined (Chapter~\ref{ch:rc-rl}) or combined with inductors
(Chapter~\ref{ch:rlc-resonance}), since it lets both magnitude and
phase be tracked with ordinary complex-number arithmetic.

\section{Capacitor Types and Practical Considerations}

\begin{itemize}
\item \textbf{Ceramic}: small, cheap, low to moderate values (pF to
low $\mu$F), widely used for RF bypass and coupling; some formulations
(especially high-K types) have capacitance that varies with applied
voltage and temperature.
\item \textbf{Film} (polyester, polypropylene): good stability and low
loss, common in filters and audio circuits.
\item \textbf{Electrolytic}: large capacitance (µF to mF) in a small
volume, but \textbf{polarized} --- must be connected with correct
polarity or they can fail catastrophically --- and generally unsuitable
for RF due to significant internal series resistance/inductance at high
frequency.
\item \textbf{Tantalum}: also polarized, more stable than aluminum
electrolytics but historically less tolerant of voltage transients.
\item \textbf{Variable (trimmer/tuning)}: mechanically adjustable
capacitance, used for tuning and neutralization.
\end{itemize}

\begin{safetynote}[Electrolytic capacitor polarity and residual charge]
Reverse-connecting a polarized capacitor, or exceeding its voltage
rating, can cause it to vent or explode. Large electrolytic and film
capacitors in power supplies can also retain a dangerous stored charge
long after power is removed; always verify a capacitor is discharged
(with an appropriate bleeder resistor or a meter check) before handling.
\end{safetynote}

\section{Real Capacitors: ESR and Tolerance}

A real capacitor has a small \textbf{equivalent series resistance (ESR)}
and series inductance, both of which become significant at high
frequencies and ultimately give every capacitor a \textbf{self-resonant
frequency} above which it behaves inductively rather than capacitively.
Capacitor tolerance is typically wider than resistor tolerance (often
$\pm10\%$ or worse for ceramic types), a factor to keep in mind when
designing or troubleshooting frequency-sensitive circuits such as filters
and oscillators.

\begin{quickreview}
\item $C = Q/U$; parallel-plate capacitance $C = \varepsilon_r
\varepsilon_0 A/d$.
\item $i = C\,dU/dt$: capacitors block DC and initially pass sudden
voltage steps as if shorted.
\item Energy stored: $W = \tfrac12 CU^2$.
\item Parallel capacitors add; series capacitors combine by reciprocal
sum (opposite of resistors).
\item $X_C = 1/(2\pi f C)$: reactance falls as frequency rises; current
leads voltage by $90\degree$; $Z_C = -jX_C$.
\end{quickreview}

\begin{practiceproblems}
\item Find the capacitance of a capacitor storing 5\,$\mu$C at 12\,V.
\item Two capacitors, 47\,nF and 100\,nF, are in parallel. Find the total
capacitance.
\item Find the reactance of a 220\,pF capacitor at 3.5\,MHz and at
28\,MHz. Comment on how reactance changes with frequency.
\item A 1000\,$\mu$F electrolytic capacitor is charged to 25\,V. Find the
stored energy, and explain why this capacitor should be treated as a
potential shock hazard even after the power supply is switched off.
\item Explain, using the $i = C\,dU/dt$ relationship, why a capacitor
appears as a short circuit to a very fast voltage transient (such as a
static discharge) even though it blocks steady DC.
\end{practiceproblems}
